\documentclass[a4paper]{article}

\usepackage[left=2cm, right=2cm, top=2cm, bottom=2cm]{geometry} 
\usepackage[utf8]{inputenc} 
\usepackage{german} 
\usepackage{amsmath} 
\usepackage{amssymb} 
\usepackage{listings} 
\usepackage{hyperref} 
\usepackage{graphicx} 
\usepackage{eurosym}
\usepackage{wrapfig} 
\usepackage{xspace}
\usepackage{ulem} \DeclareUnicodeCharacter{20AC}{\euro}

% \usepackage{helvet}
\renewcommand{\familydefault}{\sfdefault}

\newcommand{\name}{Ursula de Boor\xspace}
\newcommand{\ort}{Kirchhain\xspace}


\begin{document}
\thispagestyle{empty}
    
\begin{center}
    \begin{minipage}{.5\textwidth}
      \flushleft
      \includegraphics[width=.3\linewidth]{Regionsbadge.pdf}
    \end{minipage}%
    \begin{minipage}{.5\textwidth}
      \flushright
      \includegraphics[width=.3\linewidth]{Zeichen.pdf}
    \end{minipage}
    % ~\\[1cm]
    
    
    
	{\Huge \textbf{Stammesordnung}}\\
    % \small{Version nach der Mitgliederversammlung am 10. Mai 2015}
    ~\\[2cm]
    \includegraphics[width=0.8\textwidth]{WeisseRose.png}
    ~\\[2,5cm] 
    
    \normalsize{Verband Christlicher Pfadfinderinnen und Pfadfinder}\\
	\Huge{\textbf{Stamm \name \ort}}
\end{center}


\clearpage
\tableofcontents 
\clearpage

\begin{wrapfigure}{r}{.4\linewidth}
    \vspace{-10pt}
    \centering
    \includegraphics[width=1\linewidth]{WeisseRose.png}
    \caption{Stammeszeichen: Weiße Rose}
\end{wrapfigure}
    
\section{Aufgabe und Mitgliedschaft}
\begin{itemize}
	\item Der VCP Stamm \name \ort ist ein Teil des Verbandes Christlicher Pfadfinderinnen und Pfadfinder (VCP) mit Sitz im Wichernweg 3, 34221 Kassel. 
	\item Die Aufgaben und Ziele richten sich nach der Bundesordnung des VCP und werden gegebenen falls durch die bestehende Landesordnung des Landes Hessen sowie die Regionsordnung der Region Kurhessen ergänzt. 
	\item Mitglied des VCP Stammes \name \ort ist, wer beim Bundesverband angemeldetes Mitglied ist und dort als zum Stamm \name \ort gehörig registriert ist. Es gelten die Regelungen zu Beginn und Ende der Mitgliedschaft im VCP. 
\end{itemize}

\section{Allgemeines}
\begin{itemize}
    \item Der Stamm \name hat seinen Sitz in der Evangelischen Kirchengemeinde Kirchhain, Hinterm Kirchhof 25, 35274 Kirchhain.
    \item Unsere Pfadfinderarbeit ist offen für Menschen aller Konfessionen und Religionen bzw. nicht-religiösen Menschen. Der VCP ist ein koedukativer Verband, alle Geschlechter sind gleichberechtigt, unsere Arbeit findet gemeinsam mit allen statt.
    \item Der Stamm in Kirchhain hat sich selbst den Namen ''\name'' gegeben. \name wurde in Kirchhain geboren und war ab 1941 Mitglied der Weißen Rose Hamburg und der \textit{Candidates of Humanity}, einer Gruppe von Medizinern die in Opposition zum NS-Regime standen. 
    % \item Der Stamm \name hat sich selbst ein Zeichen gegeben. Das Stammeswappen ist nach dem Namensgeber des Stammes die Lutherrose. Sie ist wie in Abbildung \ref{lutherrose} zu verwenden.
    % \item Der VCP ist ein koedukativer Verband, in dem beide Geschlechter gleichberechtigt sind. Für die bessere Verständlichkeit wird in der Stammesordnung allerdings nicht nach Geschlechtern unterschieden.
\end{itemize}

\section{Die Gruppen des Stammes}
In dieser Ordnung werden nur strukturelle Aspekte der Arbeit erläutert, näheres zu den Arbeitsweisen der einzelnen Stufen ist in der Stufenkonzeption des VCP nachzulesen. Der Stamm nimmt die Stufenkonzeption des VCP an und versucht sie bestmöglich in der aktiven Arbeit umzusetzen. Daraus ergeben sich folgende Gruppen für den Stamm:
    \subsection{Die Kinderstufe (Wölflinge)}
		\begin{itemize}
			\item Die Kinderstufe nimmt Kinder im Alter von 7 bis 10 Jahren auf. In Ausnahmefällen kann die Aufnahme von Kindern im Alter von unter 7 Jahren durch die Führungsrunde genehmigt werden. 
            \item Mitglieder der Meute werden als Wölflinge bezeichnet. Sie tragen das rot-blaue Halstuch.
			\item Die Arbeitsweisen der Kinderstufe sind auf Rudyard Kiplings Roman „Das Dschungelbuch“ gestützt. Daraus ergeben sich auch die Bezeichnungen in dieser Stufe. 
			\item Die gesamte Kinderstufe des Stammes bezeichnet sich als Wölflinge.
			\item Der Leiter der Meute wird als Akela bezeichnet. Bei seiner Arbeit unterstützen ihn die Meutenhelfer. In den Aufgabenbereich des Akela fallen die Planung und Durchführung der Meutenstunden sowie sämtlicher anderer Aktion der Meute. 
		\end{itemize}
    \subsection{Die Pfadfinderstufe}
    \begin{itemize}
        \item Die Pfadfinderstufe ist für Kinder und Jugendliche im Alter von 10 bis 16 Jahren. 
        \item Die Arbeit erfolgt in kleinen Gruppen von 4 bis 10 Personen, die Sippen genannt werden. Alle Jungpfadfinder- und Pfadfinder*innensippen des Stammes werden als Trupp bezeichnet.
        \item Jede Sippe besitzt mindestens eine*n Sippenführer*in, der oder die für die Planung und Durchführung der Sippenstunden sowie sämtlicher anderer Aktionen der Sippe zuständig sind.
        \item Die Pfadfinderstufe ist in zwei Phasen untergliedert:
        \subsubsection{Jungpfadfinderphase}
            \begin{itemize}
                \item Die Jungpfadfinderphase nimmt Kinder im Alter von 10 bis 13 auf. Sie tragen das hellgrün-blaue Halstuch 
                \item Schwerpunkt der Arbeit liegt auf dem zusammenwachsen der Gruppe und dem spielerischen Lernen. Sie nehmen an Veranstaltungen der Region und des Landes teil.
            \end{itemize}
        \subsubsection{Pfadfinderphase}
        \begin{itemize}
            \item Die Pfadfinderphase nimmt Jugendliche im Alter von 13 bis 16 auf. Mitglieder dieser Stufe tragen das dunkelgrün-blaue Halstuch.
            \item In der Pfadfinderphase übernehmen die Jugendlichen eigene Verantwortung. Schwerpunkt ist weiterhin die kleine Gruppe und insbesondere die persönliche Entwicklung des einzelnen.
        \end{itemize}
    \end{itemize}
    
    \subsection{Ranger/Rover Stufe}
        \begin{itemize}
            \item Die Ranger/Rover Stufe nimmt Jugendliche im Alter von 16-21 Jahren auf. Sie tragen das bordeauxrot-blaue Halstuch.
            \item Die Stufe übernimmt überwiegend leitende Aufgaben im Stamm und organisiert sich selbstständig in einer oder mehreren Roverrunden. 
            \item Die Ranger/Rover erarbeiten sich eigene Projekte und nehmen überwiegend an überregionalen Veranstaltungen des VCP teil. 
            \item Die Roverrunden wählen je einen Roversprecher aus sich selbst heraus, der die Gruppe in der Führungsrunde vertritt.
        \end{itemize}
    
    \subsection{Erwachsenenstufe}
	\begin{itemize}
        \item Die Erwachsenenstufe steht allen Mitgliedern ab 21 Jahren offen. Der Übergang zur Erwachsenenstufe findet findet in persönlicher Art und Weise statt.
        \item Die Mitglieder der Erwachsenenarbeit tragen das violett-blaue Halstuch.
		\item Die Form der Erwachsenenarbeit ist je nach Interesse zu gestalten.
        \item Die Erwachsenenarbeit organisiert sich selbst und kann einen Erwachsenensprecher in die Führungsrunde entsenden.
	\end{itemize}

\section{Gremien}
    \subsection{Führungsrunde}
	\begin{itemize}
		\item Die Führungsrunde setzt sich aus dem*der Stammesführer*in, dem*der Materialwart*wärtin, den aktiven Gruppenleiter*innen und den Sprecher*innen der Ranger/Rover und der Erwachsenen zusammen. 
		\item Freie Mitarbeiter*innen und Meutenhelfer*innen können von der Führungsrunde eingeladen werden. 
		\item Die Führungsrunde wird von dem*der Stammesführer*in einberufen und sollte mindestens einmal im Monat stattfinden. 
		\item Die Führungsrunde versucht nach Möglichkeit ohne Abstimmung Einigungen zu finden. Sollte es dennoch zu einen Abstimmungsverfahren kommen, so gilt die einfache Mehrheit ohne Berücksichtigung der Enthaltungen als beschlussfähig. 
		\item Die Führungsrunde hat das Recht, Punkte der Stammesordnung außer Kraft zu setzen. 
		\item Des weiteren unterliegt die Führungsrunde einer Schweigepflicht und ist somit auch für vertrauliche Gespräche zugänglich. Daraus ergibt sich das Recht der Führungsrunde, Personen aus der Debatte zu bestimmten Themen auszuschließen. 
		\item Die Protokolle der Führungsrunde sind nur für Mitglieder*innen der Führungsrunde einsehbar. Die Führungsrunde kann das Protokoll auf Anfrage zur Einsicht freigeben.
		\subsubsection{Aufgaben der FüRu} 
		\begin{itemize}
		    \item Aufstellen eines Jahresplanes 
			\item Vorbereitung der Stammesaktivitäten 
			\item Öffentlichkeitsarbeit 
			\item Innerverbandliche Informationsvermittlung 
			\item Beauftragen von Vertretern für Gremien der Region und des Landes 
            \item Langfristige Planung des Stammes
		\end{itemize}
	\end{itemize}
    
    \subsection{Stammesführer*in}
	\begin{itemize}
		\item Der*die Stammesführer*in muss volljährig und Mitglied im VCP sein. 
		\item Die Amtszeit beträgt 2 Jahre. Wahl und Entlastung findet durch die Mitgliederversammlung statt.
		\item Der*die Stammesführer*in hat den Vorstand über die Führungsrunde. Er beruft die Führungsrunden, moderiert sie und hat für die Erstellung eines Protokolls Sorge zu tragen. 
		\item Der*die Stammesführer*in ist ermächtigt, zwischen den Führungsrunden weitreichende Entscheidungen den Stamm betreffend eigenmächtig zu entscheiden, muss sich für diese aber der Mitgliederversammlung und der Führungsrunde gegenüber verantworten. 
		\item Der Stammesführer hat die Pflicht an den Tagungen des Regionsrates teilzunehmen. 
		\item Kann der*die Stammesführer*in bestimmte Aufgaben nicht wahrnehmen, ist es möglich diese an Mitglieder der Führungsrunde und freie Mitarbeiter zu delegieren. 
	\end{itemize}
    
    \subsubsection{Materialwart}
	\begin{itemize}
		\item Der*die Materialwart*wärtin muss nicht volljährig, aber Mitglied im VCP sein. 
		\item Der*die Materialwart*wärtin verwaltet das Material. 
        \item Die Amtszeit beträgt 2 Jahre. Wahl und Entlastung findet durch die Mitgliederversammlung statt.
		\item Eine Ausgabe von Material für privaten Gebrauch kann nur mit Genehmigung der Führungsrunde erfolgen. Der*die Abholer*in haftet mit seiner Unterschrift für das Material. 
		\item Der*die Materialwart*wärtin gibt der Mitgliederversammlung jährlich einen Bericht. 
	\end{itemize}
    
    \subsection{Freie Mitarbeiter}
	\begin{itemize}
		\item Jeder Ranger/Rover des Stammes ohne festes Amt ist als freier Mitarbeiter zu sehen. 
		\item Freie Mitarbeiter können bestimmte Aufgaben wahrnehmen, die durch die Führungsrunde beschlossen werden und sich auf Regions- und Landesebene engagieren. 
	\end{itemize}

\section{Mitgliederversammlung}

    \subsection{Einberufung und Beschlussfähigkeit}
	\begin{itemize}
        \item Die Mitgliederversammlung tritt regulär einmal im Jahr zusammen und wird durch den*die Stammesführer*in einberufen und moderiert.
        \item Weitere Mitgliederversammlungen werden auf Verlangen der Führungsrunde oder der einfachen Mehrheit der eingetragenen Mitglieder*innen einberufen.
        \item Die Mitgliederversammlung setzt sich aus den angemeldeten (Jung-)Pfadfinder*innen, Ranger/Rover und Erwachsenen des Stammes zusammen.
        \item Die Mitgliederversammlung ist beschlussfähig, wenn sie ordnungsgemäß einberufen wurde.
        \item Die Einberufung der Mitgliederversammlung erfolgt mindestens 14 Tage vor dem Versammlungstermin durch schriftliche oder elektronische Einladung an jedes Mitglied.
        \item Wahlen und Beschlüsse erfolgen mit der einfachen Mehrheit ohne Berücksichtigung der Stimmenthaltungen.
        \item Änderungen der Stammesordnung benötigen eine 2/3 Mehrheit der Anwesenden der Mitgliederversammlung ohne Berücksichtigung der Stimmenthaltungen.
        \item Die Mitgliederversammlung tagt öffentlich, somit ist es jedem möglich an der Mitgliederversammlung teilzunehmen und Einblick in die Stammesarbeit zu bekommen.
        \item Jeder ist berechtigt, Anträge an die Mitgliederversammlung zu stellen.
	\end{itemize}
    
    \subsection{Aufgaben} 
	\begin{itemize}
        \item Die Mitgliederversammlung wählt eine*n Stammesführer*in, eine*n Materialwart*wärtin und weitere Stammesbeauftragte.
        \item Die Mitgliederversammlung nimmt den Bericht des*der Stammesführer*in und der Führungsrunde entgegen. Es wird über das vergangene Jahr und die kommende Jahresplanung, die Gruppen des Stammes, die Anbindung an die Kirchengemeinde, sowie die Finanzen des Stammes berichtet.
        \item Die Sprecher*innen der Ranger/Rover und der Erwachsenenarbeit berichten über ihre Arbeit.
        \item Die Mitgliederversammlung kann Entlastungen erteilen.
	\end{itemize}

\clearpage
\begin{appendix}
\section{Anhang}

\subsection{Versionshistorie}
\subsubsection{Mitgliederversammlung am 13. Dezember 2019}
\begin{itemize}
    \item Beschluss der Stammesordnung in der ersten Fassung.
\end{itemize}
    

\end{appendix}






















\end{document}
